\section{Justificación}
{\justify
Actualmente, el proceso de orientación para estudiantes nuevos, visitantes y padres de familia dentro del campus universitario de la Universidad del Valle de Guatemala presenta diversas limitaciones. Tradicionalmente, el recorrido por las instalaciones es guiado por estudiantes o personal universitario, pero este se realiza una única vez —si es que se llega a realizar— lo cual suele ser insuficiente para que los usuarios recuerden rutas, espacios o servicios disponibles. Como consecuencia, durante los primeros semestres, muchos estudiantes enfrentan dificultades para ubicarse, llegando a perderse o desaprovechar recursos valiosos del campus.

Además de la desorientación inicial, existen barreras de accesibilidad que afectan a personas con movilidad reducida o con discapacidades visuales, quienes podrían beneficiarse significativamente de un sistema que les permita planificar sus trayectos de forma precisa y adaptada a sus necesidades. Una solución tecnológica puede contribuir a solventar ambos problemas, mejorando tanto la orientación como la inclusión dentro del entorno universitario.

En este contexto, el desarrollo de una aplicación de realidad aumentada con capacidades de localización y navegación precisa representa una oportunidad para aplicar conocimientos técnicos adquiridos a lo largo de la carrera de Ingeniería en Ciencias de la Computación y Tecnologías de la Información. Este proyecto integra áreas clave como programación orientada a objetos, interacción humano-computadora, gráficos por computadora, desarrollo móvil y algoritmos de búsqueda de caminos (pathfinding). Asimismo, se utilizarán sensores Estimote Beacons con tecnología Ultra-Wideband (UWB), cuya precisión y facilidad de conexión brindan una ventaja considerable en comparación con tecnologías previamente utilizadas, como los códigos QR.

Este trabajo, al ser parte de un megaproyecto en curso, busca fortalecer y mejorar la versión previa de la aplicación desarrollada por generaciones anteriores, proponiendo avances significativos en la precisión de la localización y en la experiencia del usuario. La solución no solo resolverá problemas existentes, sino que también funcionará como base tecnológica para futuras generaciones de estudiantes que podrán continuar desarrollando y afinando la herramienta.

Finalmente, este proyecto aporta a la universidad al fortalecer su imagen como institución innovadora, capaz de generar soluciones tecnológicas útiles y vanguardistas dentro de su propio entorno. Comunica el potencial de sus estudiantes para abordar problemas reales mediante el uso creativo y riguroso del conocimiento adquirido, posicionando a la UVG como una universidad comprometida con la tecnología, la accesibilidad y la experiencia estudiantil.}

\newpage